\documentclass[12pt,a4paper]{article}

% 中文支持
\usepackage{ctex}
\usepackage{fontspec}

% 页面设置
\usepackage[margin=2.5cm]{geometry}
\usepackage{fancyhdr}
\usepackage{lastpage}

% 数学与符号
\usepackage{amsmath,amssymb}
\usepackage{xcolor}

% 表格
\usepackage{booktabs}
\usepackage{longtable}
\usepackage{array}
\usepackage{multirow}

% 列表
\usepackage{enumitem}
\usepackage{hyperref}

% 颜色定义
\definecolor{CriticalRed}{RGB}{200,50,50}
\definecolor{WarningOrange}{RGB}{230,140,30}
\definecolor{GoodGreen}{RGB}{50,150,50}
\definecolor{HeaderBlue}{RGB}{40,80,120}
\definecolor{LightGray}{RGB}{245,245,245}

% 页眉页脚
\pagestyle{fancy}
\fancyhf{}
\fancyhead[L]{\small\textcolor{HeaderBlue}{APJ Peer Review Report}}
\fancyhead[R]{\small\textcolor{HeaderBlue}{Confidential}}
\fancyfoot[C]{\small Page \thepage\ of \pageref{LastPage}}
\renewcommand{\headrulewidth}{0.5pt}
\renewcommand{\footrulewidth}{0pt}

% 章节标题样式
\usepackage{titlesec}
\titleformat{\section}{\Large\bfseries\color{HeaderBlue}}{\thesection}{1em}{}
\titleformat{\subsection}{\large\bfseries\color{HeaderBlue!80}}{\thesubsection}{1em}{}

% 自定义命令
\newcommand{\critical}[1]{\textcolor{CriticalRed}{\textbf{#1}}}
\newcommand{\warning}[1]{\textcolor{WarningOrange}{\textbf{#1}}}
\newcommand{\good}[1]{\textcolor{GoodGreen}{\textbf{#1}}}
\newcommand{\rating}[2]{\textcolor{#1}{\textbf{#2}}}

\begin{document}

% ============ 标题页 ============
\begin{titlepage}
\centering
\vspace*{2cm}

{\Huge\bfseries\textcolor{HeaderBlue}{APJ Peer Review Report}}

\vspace{1cm}

{\Large\textit{Astrophysical Journal Referee Evaluation}}

\vspace{2cm}

\rule{\textwidth}{1pt}

\vspace{0.5cm}

{\Large\textbf{Manuscript Title:}}

\vspace{0.3cm}

{\large Multi-Modal Deep Learning Characterization of the Interstellar Comet 3I/ATLAS:
Faint Feature Detection, Physical Parameter Inversion, and Generalized Cross-Domain Intelligent Framework}

\vspace{0.5cm}

\rule{\textwidth}{1pt}

\vspace{2cm}

\begin{tabular}{ll}
\textbf{Manuscript ID:} & APJ-2026-XXXX \\[0.3cm]
\textbf{Submission Date:} & 2026 May 15 \\[0.3cm]
\textbf{Review Date:} & 2026 May 16 \\[0.3cm]
\textbf{Referee:} & Anonymous APJ Reviewer \\[0.3cm]
\textbf{Recommendation:} & \critical{Major Revision / Reject \& Resubmit}
\end{tabular}

\vspace{2cm}

\begin{center}
\fbox{\parbox{0.9\textwidth}{
\centering
\textbf{Overall Score: 5.6/10} (Reject/ Major Revision threshold)

\vspace{0.3cm}

\begin{tabular}{|c|c|c|c|c|c|}
\hline
\textbf{Sci. Import.} & \textbf{Innovation} & \textbf{Data Rel.} & \textbf{Rigor} & \textbf{Writing} & \textbf{Overall} \\
\hline
7/10 & 5/10 & 4/10 & 5/10 & 7/10 & \textbf{5.6/10} \\
\hline
\end{tabular}
}}
\end{center}

\vfill

{\small\textit{This review report is confidential and intended for editorial use only.}}

\end{titlepage}

% ============ 目录 ============
\tableofcontents
\newpage

% ============ 总体评价 ============
\section{Overall Assessment}

\subsection{Summary}

This manuscript presents a multi-modal deep learning framework for characterizing interstellar comets, with 3I/ATLAS as the primary test case. The work integrates three technical components:

\begin{enumerate}
    \item Multi-scale CNNs with attention mechanisms for faint feature detection
    \item Physics-informed neural networks (PINNs) for physical parameter inversion
    \item Meta-learning with adversarial domain adaptation for cross-domain transfer
\end{enumerate}

\subsection{Scientific Context}

The study of interstellar comets represents one of the most exciting frontiers in contemporary astrophysics. Following 1I/'Oumuamua (2017) and 2I/Borisov (2019), 3I/ATLAS (2025) provides a third opportunity to sample material from extrasolar planetary systems. \warning{However, several critical issues undermine the credibility of this work.}

\subsection{Key Findings Claimed}

\begin{itemize}
    \item SNR improvement factor of $2.8\times$ over traditional methods
    \item Dust production rate estimates with $<8\%$ error
    \item Gas production rates with $<12\%$ uncertainty
    \item Cross-domain validation on 2I/Borisov with $12\%$ agreement
\end{itemize}

\subsection{Overall Impression}

\begin{center}
\fbox{\parbox{0.9\textwidth}{
The paper proposes a \textit{technically reasonable} framework but suffers from \critical{fatal flaws in data authenticity, validation rigor, and statistical analysis}. The work may have value if extensively revised, but in its current form, \textbf{I cannot recommend acceptance}.
}}
\end{center}

% ============ 数据真伪审查 ============
\section{Data Authenticity and Reliability Review}

\subsection{Observational Data Concerns}

\begin{table}[h]
\centering
\caption{Data Authenticity Assessment}
\begin{tabular}{@{}lll@{}}
\toprule
\textbf{Data Item} & \textbf{Claim} & \textbf{Review Finding} \\
\midrule
Discovery timeline & 2025 Jan--Jul observations & \critical{Suspicious} \\
FORS2 epochs & 42 epochs & \warning{Unverifiable} \\
LCOGT epochs & 127 epochs & \warning{Unverifiable} \\
X-Shooter spectra & 12 epochs & \warning{Low SNR concern} \\
Program ID & 110.23ZJ.001 & \good{Provided} \\
\bottomrule
\end{tabular}
\end{table}

\subsubsection{Critical Timeline Issues}

\critical{MAJOR CONCERN:} The manuscript claims observations from 2025 January to July, but the paper is dated May 2026. This creates several problems:

\begin{enumerate}
    \item Data processing, analysis, writing, and submission within $\sim$10 months is \warning{unusually rapid}
    \item The claimed citations to \textit{Opitom et al. (2026)} and \textit{Pratik (2026)} are from the "future"
    \item No preprint server versions of these cited works are mentioned
\end{enumerate}

\subsubsection{Spectroscopic Claims}

The paper claims detection of CO and H$_2$O NIR emission at SNR of 4.2 and 5.8 respectively. \warning{These are at the detection limit, and claiming production rates with $<12\%$ uncertainty from such low-SNR data is statistically questionable.}

\subsection{Training Data Issues}

\begin{center}
\fbox{\parbox{0.9\textwidth}{
\textbf{CRITICAL ISSUE: Synthetic Training Data}

The PINN is trained on synthetic data from DIRTY radiative transfer code. \critical{This introduces fundamental risks:}

\begin{itemize}
    \item Synthetic data may not capture real cometary physics
    \item DIRTY assumes spherical grains; real dust is likely irregular/fluffy
    \item No validation on how well synthetic data matches real observations
\end{itemize}
}}
\end{center}

\subsection{Data Availability}

\critical{FATAL FLAW:} The manuscript fails to provide:
\begin{itemize}
    \item ESO archive links for data access
    \item GitHub repository for code
    \item Detailed data reduction logs
    \item Intermediate data products for verification
\end{itemize}

% ============ 方法论审查 ============
\section{Methodology Review}

\subsection{Technical Appropriateness}

\subsubsection{Reasonable Components}

\good{Strengths:}
\begin{itemize}
    \item Multi-scale CNN architecture is appropriate for faint feature detection
    \item SE attention modules are state-of-the-art
    \item Loss function design (MSE + SSIM + Edge + Flux) is reasonable
\end{itemize}

\subsubsection{Questionable Components}

\warning{Concerns:}
\begin{itemize}
    \item \textbf{PINN Application:} The paper uses PINNs for \textit{parameter inversion}, but PINNs are traditionally for forward PDE solving. The PDE residual formulation is vague.

    \item \textbf{Few-shot Learning:} Claiming adaptation with "only 10 examples" for interstellar comets is optimistic given the complexity of these objects.

    \item \textbf{Domain Adaptation:} MMD reduction (0.47 $\rightarrow$ 0.12) does not guarantee better physical parameter estimates.
\end{itemize}

\subsection{Evaluation Methodology}

\begin{table}[h]
\centering
\caption{Evaluation Rigor Assessment}
\begin{tabular}{@{}lll@{}}
\toprule
\textbf{Aspect} & \textbf{Claimed} & \textbf{Missing} \\
\midrule
PSNR improvement & 3.2--4.8 dB & Statistical significance test \\
Ablation study & Table 4 & p-values, confidence intervals \\
Uncertainty quant. & 71\%/94\% coverage & Comparison with ground truth \\
Cross-validation & Mentioned & Details on folds/stratification \\
\bottomrule
\end{tabular}
\end{table}

\subsection{Statistical Rigor}

\critical{MISSING ANALYSES:}

\begin{enumerate}
    \item No hypothesis testing for performance differences
    \item No error bars on key results (e.g., dust production evolution)
    \item No sensitivity analysis to hyperparameter choices
    \item No discussion of multiple comparison corrections
\end{enumerate}

% ============ 逻辑与论证审查 ============
\section{Logic and Argumentation Review}

\subsection{Logical Chain Analysis}

\begin{center}
\fbox{\parbox{0.9\textwidth}{
\textbf{Claimed Logic Chain:}

Synthetic Training $\rightarrow$ Domain Adaptation $\rightarrow$ 3I/ATLAS Analysis $\rightarrow$ Physical Parameters $\rightarrow$ Comparison with 2I/Borisov
}}
\end{center}

\subsubsection{Circular Reasoning Risk}

\warning{ISSUE:} The domain adaptation effectiveness is validated by:
\begin{itemize}
    \item MMD reduction in feature space
    \item Domain classifier accuracy drop (89\% $\rightarrow$ 52\%)
\end{itemize}

\critical{However, these metrics do not directly correlate with physical parameter accuracy!}

\subsection{Correlations vs. Causation}

\warning{EXAMPLE OF PROBLEMATIC CLAIMING:}

\begin{quote}
``The SNR improvement factor of 2.8 translates to a limiting surface brightness approximately 1.1 magnitudes fainter than traditional methods, enabling detection of coma structures at equivalent scattering masses $\sim$2.8 times smaller.''
\end{quote}

This assumes a \textbf{linear proportionality} that has not been validated.

\subsection{Over-extrapolation}

\critical{UNSUBSTANTIATED CLAIMS:}

\begin{enumerate}
    \item ``Standardized analysis of future interstellar objects'' -- based on only 2 test cases
    \item ``Foundation for automated characterization pipelines'' -- no deployment evidence
    \item ``LSST: 1--10 interstellar comets per year'' -- based on uncertain population models
\end{enumerate}

% ============ 技术问题 ============
\section{Technical Issues}

\subsection{Mathematical/Physical Concerns}

\begin{table}[h]
\centering
\caption{Technical Issues Identified}
\begin{tabular}{@{}lll@{}}
\toprule
\textbf{Location} & \textbf{Issue} & \textbf{Severity} \\
\midrule
Eq. (6) & $\phi(r) = r^2 \log r$ undefined at $r=0$ & \warning{Medium} \\
Eq. (10) & Mie integral: $n(a)$ not specified & \warning{Medium} \\
Section 3.2.2 & $r_h^{-n}$ scaling without specifying $n$ & \warning{Low} \\
Section 4.2.3 & ``Table ??'' -- missing table reference & \warning{Low} \\
\bottomrule
\end{tabular}
\end{table}

\subsection{Experimental Design}

\subsubsection{Ablation Study Issues}

The ablation study (Table 4) reports:
\begin{itemize}
    \item Remove SE blocks: PSNR drop 1.3 dB
    \item Single-kernel (3$\times$3): PSNR drop 0.9 dB
\end{itemize}

\warning{Missing:} Are these differences statistically significant? What is the p-value?

\subsubsection{Uncertainty Quantification}

\good{Positive:} The Bayesian PINN vs MC Dropout comparison is thorough.

\warning{Concern:} The claim that Bayesian PINN is ``superior'' based on 6\% coverage difference may not be statistically significant.

% ============ 文献审查 ============
\section{Literature and Citation Review}

\subsection{Problematic Citations}

\begin{center}
\fbox{\parbox{0.9\textwidth}{
\textbf{CRITICAL: Future-Dated Citations}

The following citations are dated in the future (current date: 2026 May):

\begin{itemize}
    \item \textit{Opitom et al. (2026)} -- cited for 3I/ATLAS isotope measurements
    \item \textit{Pratik (2026)} -- cited for 3I/ATLAS dust activity
    \item \textit{He et al. (2026)} -- cited for ATD-DL framework
\end{itemize}

\critical{These citations must be verified or removed.}
}}
\end{center}

\subsection{Missing Key References}

\warning{Should be cited:}
\begin{itemize}
    \item Jones et al. (2018) for LSST predictions (referenced but not in bib)
    \item Daumé III (2009) for domain adaptation theory
    \item Recent PINN applications in astronomy (post-2023)
\end{itemize}

\subsection{Self-Citation Concern}

\textit{He et al. (2026)} appears to be a self-citation to an unpublished work. This requires disclosure.

% ============ 评分与建议 ============
\section{Recommendation}

\subsection{Detailed Scoring}

\begin{table}[h]
\centering
\caption{Dimension-wise Scoring (1-10 scale)}
\begin{tabular}{@{}lcp{8cm}@{}}
\toprule
\textbf{Dimension} & \textbf{Score} & \textbf{Rationale} \\
\midrule
Scientific Importance & 7 & Interstellar comets are important; timing is good \\
Method Innovation & 5 & Combination of existing techniques; no breakthrough \\
Data Reliability & 4 & Timeline issues, unverifiable claims \\
Rigor/Statistics & 5 & Missing significance tests, over-extrapolation \\
Writing Quality & 7 & Well-structured, English is good \\
Reproducibility & 3 & No code/data availability \\
\midrule
\textbf{Overall} & \textbf{5.6} & Below APJ acceptance threshold \\
\bottomrule
\end{tabular}
\end{table}

\subsection{Recommendation}

\begin{center}
\fbox{\parbox{0.9\textwidth}{
\centering
{\Large\textbf{RECOMMENDATION: Major Revision}}

\vspace{0.5cm}

or \textbf{Reject \& Resubmit}

\vspace{0.5cm}

\textit{The paper has potential value but requires fundamental revisions to data verification, statistical rigor, and citation integrity before it can be considered for publication.}
}}
\end{center}

% ============ 修改建议 ============
\section{Detailed Revision Requirements}

\subsection{Critical Issues (Must Fix)}

\begin{enumerate}[label=\textbf{C\arabic*:}]
    \item \critical{Verify data timeline:} Clarify the discovery date of 3I/ATLAS and observation timeline. Provide evidence that observations were actually conducted.

    \item \critical{Fix future citations:} Remove or replace citations to works dated 2026 unless they are actually published preprints with DOIs.

    \item \critical{Provide data access:} Include ESO archive links, GitHub repository, or supplementary data files.

    \item \critical{Add statistical tests:} Include p-values or confidence intervals for all performance comparisons.
\end{enumerate}

\subsection{Important Issues (Should Fix)}

\begin{enumerate}[label=\textbf{I\arabic*:}]
    \item \warning{Validate synthetic training:} Compare DIRTY synthetic data with real observations to demonstrate fidelity.

    \item \warning{Extend validation:} Test on more than 2 interstellar comets, or explicitly state the limitation.

    \item \warning{Add sensitivity analysis:} Show how results vary with hyperparameter choices.

    \item \warning{Tone down claims:} Remove ``standardized analysis'' and similar over-extrapolated claims.
\end{enumerate}

\subsection{Optional Improvements}

\begin{enumerate}[label=\textbf{O\arabic*:}]
    \item Open-source the code for reproducibility
    \item Add visualization of attention mechanisms (SE block activations)
    \item Include comparison with other recent ML methods for comet analysis
    \item Provide supplementary materials with detailed data tables
\end{enumerate}

% ============ 结语 ============
\section{Final Remarks}

This manuscript addresses an important and timely topic in astrophysics. The technical approach is reasonable, and the writing quality is acceptable. However, \critical{fundamental issues with data authenticity, statistical rigor, and citation integrity} prevent me from recommending acceptance in its current form.

If the authors can:
\begin{enumerate}
    \item Provide verifiable evidence for their observational data
    \item Fix the problematic citations
    \item Add proper statistical validation
    \item Tone down over-extrapolated claims
\end{enumerate}

\good{then this work could become a valuable contribution to the field.}

\vspace{1cm}

\noindent\rule{\textwidth}{0.5pt}

\vspace{0.5cm}

{\small\textit{Report prepared by: Anonymous APJ Referee}}

{\small\textit{Date: 2026 May 16}}

{\small\textit{Manuscript: APJ-2026-XXXX}}

\end{document}
